\documentclass[12pt]{article}
% Préambule commun à tous les documents extraits de « La Revue CAPES »
\usepackage[T1]{fontenc}
\usepackage[utf8]{inputenc}
\usepackage{lmodern}
\usepackage{amsmath,amssymb,amsthm,amsfonts,mathrsfs}
\usepackage{setspace}
\usepackage{listings}
\usepackage{pgfplots}
\pgfplotsset{compat=1.18}
\usepackage{longtable}
\usepackage{adjustbox}
\usepackage{lettrine}
\usepackage{tkz-tab}
\usepackage{changepage}
\usepackage{pdflscape}
\usepackage{graphicx}
\usepackage{tikz}
\usepackage{xcolor}
\usepackage[a4paper,top=2cm,bottom=2.5cm,left=2.5cm,right=2cm]{geometry}
\usepackage{fancyhdr}
\usepackage{draftwatermark}
\SetWatermarkText{La RevueCapes}
\SetWatermarkLightness{0.9}
\SetWatermarkScale{2}
\usepackage{hyperref}
\hypersetup{colorlinks=true,linkcolor=blue!50!black,urlcolor=blue!50!black}

\definecolor{revue}{RGB}{20,60,120}

\pagestyle{fancy}
\fancyhf{}
\renewcommand{\headrulewidth}{0.4pt}
\setlength{\headheight}{15pt}

% \entete{Matière}{Titre}{Type de document}
\newcommand{\entete}[3]{%
  \fancyhead[L]{\small\textsc{La Revue CAPES} -- #1}%
  \fancyhead[R]{\small #2}%
  \fancyfoot[C]{\thepage}%
  \thispagestyle{plain}%
  \begin{center}
    {\color{revue}\rule{\textwidth}{1.2pt}}\\[4pt]
    {\small\textsc{Concours directs CAP-CEG / CAPES -- Burkina Faso}}\\[6pt]
    {\LARGE\bfseries\color{revue} #2}\\[6pt]
    {\large\itshape #1 -- #3}\\[2pt]
    {\color{revue}\rule{\textwidth}{1.2pt}}
  \end{center}
  \vspace{0.5cm}%
}

% Encadré pour les remarques ajoutées dans les corrigés
\newcommand{\remarque}[1]{\par\medskip\noindent\fcolorbox{revue}{revue!6}{\parbox{\dimexpr\linewidth-2\fboxsep-2\fboxrule}{\textbf{\color{revue}Remarque.} #1}}\par\medskip}


\begin{document}
\entete{Mathématiques}{Corrigé détaillé du sujet type n°4}{Correction}
\begin{spacing}{1.5}

\remarque{L'énoncé original indiquait $|a|=1$ ; il faut lire $|a|\neq1$. En effet, si $|a|=1$, alors $1+\bar a z=\bar a(a+z)$ et $f_a$ serait constante égale à $1/\bar a=a$ (et non définie en $-a$) ; de plus les valeurs $a=2$, $2i$, $1+i$ de la question 6 sont de module différent de $1$.}

\textbf{Exercice 1}

\textbf{1. $f_a$ est bien définie sur $U$.}

$f_a(z)$ est définie si et seulement si $1+\bar a z\neq0$, c'est-à-dire $\bar a z\neq-1$. Or, pour $z\in U$ :
\[
|\bar a z|=|\bar a|\,|z|=|a|\times1=|a|\neq1=|-1| .
\]
Les nombres $\bar a z$ et $-1$ n'ont pas le même module, donc $\bar az\neq-1$ : $f_a(z)$ est bien définie pour tout $z\in U$.

\textbf{2. Si $z\in U$, alors $\bar z=\dfrac1z$.}

Pour tout complexe $z$, $z\bar z=|z|^2$. Si $z\in U$, $|z|=1$ donc $z\neq0$ et $z\bar z=1$, d'où $\bar z=\dfrac1z$.

(Variante : $z=\cos\theta+i\sin\theta$, $\bar z=\cos\theta-i\sin\theta$ et $z\bar z=\cos^2\theta+\sin^2\theta=1$.)

\textbf{3. Si $z\in U$, alors $f_a(z)\in U$.}

Pour $z\in U$, on écrit $1=z\bar z$ :
\[
1+\bar az=z\bar z+\bar a z=z(\bar z+\bar a)=z\,\overline{(z+a)} .
\]
Donc
\[
f_a(z)=\frac{z+a}{z\,\overline{(z+a)}} .
\]
Remarquons que $z+a\neq0$ (sinon $a=-z$ serait de module $1$). Comme $|\overline{w}|=|w|$ pour tout complexe $w$ et que le module d'un quotient est le quotient des modules :
\[
|f_a(z)|=\frac{|z+a|}{|z|\,|z+a|}=\frac1{|z|}=1 .
\]
Donc $f_a(z)\in U$.

\textbf{4. Tout $t\in U$ a un unique antécédent dans $U$.}

Soit $t\in U$. Pour $z\in U$ :
\[
t=f_a(z)\iff t(1+\bar az)=z+a\iff t+\bar a tz=z+a\iff z(1-\bar a t)=t-a .
\]
Comme $|\bar a t|=|a|\neq1$, on a $1-\bar at\neq0$ et donc
\[
t=f_a(z)\iff z=\frac{t-a}{1-\bar at}.
\]
Ce nombre $z$ est unique. De plus, $\dfrac{t-a}{1-\bar a t}=\dfrac{t+(-a)}{1+\overline{(-a)}\,t}=f_{-a}(t)$, et comme $|-a|=|a|\neq1$, la question 3 (appliquée à $-a$) montre que $z=f_{-a}(t)\in U$.

\textbf{5. $f_a$ est une bijection de $U$ sur $U$.}

D'après les questions 3 et 4, $f_a$ envoie $U$ dans $U$ et tout élément $t$ de $U$ possède un unique antécédent $z\in U$. Donc $f_a:U\to U$ est bijective et, pour $z,t\in U$ :
\[
t=f_a(z)\iff z=f_{-a}(t).
\]
Sa bijection réciproque est $\boxed{f_a^{-1}=f_{-a}:\ t\mapsto\dfrac{t-a}{1-\bar at}}$.

\textbf{6. Antécédents des éléments de $\{-1,1,i,-i\}$.}

Les points cherchés sont les $z=f_{-a}(t)=\dfrac{t-a}{1-\bar a t}$ pour $t\in\{-1,1,i,-i\}$. On calcule en multipliant par le conjugué du dénominateur.

\underline{(a) $a=2$} : $z=\dfrac{t-2}{1-2t}$.
\begin{itemize}
\item $t=-1$ : $z=\dfrac{-3}{3}=-1$ ;\quad $t=1$ : $z=\dfrac{-1}{-1}=1$ ;
\item $t=i$ : $z=\dfrac{-2+i}{1-2i}=\dfrac{(-2+i)(1+2i)}{5}=\dfrac{-2-4i+i-2}{5}=\dfrac{-4-3i}{5}$ ;
\item $t=-i$ : $z=\dfrac{-2-i}{1+2i}=\dfrac{(-2-i)(1-2i)}{5}=\dfrac{-4+3i}{5}$.
\end{itemize}
Ensemble cherché : $\Big\{-1\,;\,1\,;\,\dfrac{-4-3i}{5}\,;\,\dfrac{-4+3i}{5}\Big\}$.

\underline{(b) $a=2i$} ($\bar a=-2i$) : $z=\dfrac{t-2i}{1+2it}$.
\begin{itemize}
\item $t=-1$ : $z=\dfrac{-1-2i}{1-2i}=\dfrac{(-1-2i)(1+2i)}{5}=\dfrac{-1-2i-2i+4}{5}=\dfrac{3-4i}{5}$ ;
\item $t=1$ : $z=\dfrac{1-2i}{1+2i}=\dfrac{(1-2i)^2}{5}=\dfrac{-3-4i}{5}$ ;
\item $t=i$ : $z=\dfrac{-i}{1-2}=i$ ;\quad $t=-i$ : $z=\dfrac{-3i}{3}=-i$.
\end{itemize}
Ensemble cherché : $\Big\{\dfrac{3-4i}{5}\,;\,\dfrac{-3-4i}{5}\,;\,i\,;\,-i\Big\}$.

\underline{(c) $a=1+i$} ($\bar a=1-i$) : $z=\dfrac{t-1-i}{1-(1-i)t}$.
\begin{itemize}
\item $t=-1$ : $z=\dfrac{-2-i}{2-i}=\dfrac{(-2-i)(2+i)}{5}=\dfrac{-4-4i+1}{5}=\dfrac{-3-4i}{5}$ ;
\item $t=1$ : $z=\dfrac{-i}{i}=-1$ ;
\item $t=i$ : le dénominateur vaut $1-(1-i)i=1-i+i^2=-i$ et le numérateur $i-1-i=-1$, donc $z=\dfrac{-1}{-i}=\dfrac1i=-i$ ;
\item $t=-i$ : $1+(1-i)i=2+i$, donc $z=\dfrac{-1-2i}{2+i}=\dfrac{(-1-2i)(2-i)}{5}=\dfrac{-4-3i}{5}$.
\end{itemize}
Ensemble cherché : $\Big\{\dfrac{-3-4i}{5}\,;\,-1\,;\,-i\,;\,\dfrac{-4-3i}{5}\Big\}$.

On vérifie que tous ces points sont bien de module $1$ (par exemple $|{-4-3i}|=5$), conformément à la question 4.

\medskip\textbf{Exercice 2}

$q_a(x_1,x_2,x_3)=x_1^2+(1+a)x_2^2+(1+a+a^2)x_3^2+2x_1x_2-2ax_2x_3$.

\textbf{1. Matrice de $q_a$ dans la base canonique.}

Les coefficients diagonaux sont les coefficients des carrés ; le coefficient $(i,j)$, $i\neq j$, est la \emph{moitié} du coefficient de $x_ix_j$ :
\[
A=\begin{pmatrix}
1 & 1 & 0 \\
1 & 1+a & -a \\
0 & -a & 1+a+a^2
\end{pmatrix},\qquad q_a(X)={}^tXAX .
\]

\textbf{2. Forme bilinéaire symétrique (forme polaire) associée.}

Pour $X=(x_1,x_2,x_3)$ et $Y=(y_1,y_2,y_3)$, $\varphi_a(X,Y)={}^tXAY$, soit
\[
\varphi_a(X,Y)=x_1y_1+(1+a)x_2y_2+(1+a+a^2)x_3y_3+x_1y_2+x_2y_1-a\,x_2y_3-a\,x_3y_2 .
\]
On vérifie que $\varphi_a(X,X)=q_a(X)$ et $\varphi_a(X,Y)=\varphi_a(Y,X)$. (On peut aussi utiliser $\varphi_a(X,Y)=\frac12\big[q_a(X+Y)-q_a(X)-q_a(Y)\big]$.)

\textbf{3. Décomposition de Gauss.}

On regroupe d'abord tous les termes contenant $x_1$, en utilisant $x_1^2+2x_1x_2=(x_1+x_2)^2-x_2^2$ :
\begin{align*}
q_a(X)&=(x_1+x_2)^2-x_2^2+(1+a)x_2^2-2ax_2x_3+(1+a+a^2)x_3^2\\
&=(x_1+x_2)^2+a\,x_2^2-2a\,x_2x_3+(1+a+a^2)x_3^2 .
\end{align*}
Puis on regroupe les termes en $x_2$ : $a(x_2^2-2x_2x_3)=a(x_2-x_3)^2-ax_3^2$ :
\[
\boxed{q_a(X)=(x_1+x_2)^2+a\,(x_2-x_3)^2+(1+a^2)\,x_3^2}.
\]
Les formes linéaires $\ell_1=x_1+x_2$, $\ell_2=x_2-x_3$, $\ell_3=x_3$ sont linéairement indépendantes (matrice triangulaire de déterminant $1$).

\textbf{4. Noyau, rang et signature.}

On rappelle que le noyau de $q_a$ est le noyau de sa forme polaire : $\ker q_a=\{X\in\mathbb{R}^3 \ /\ AX=0\}$ ; le rang de $q_a$ est le rang de $A$, et la signature $(p,q)$ compte les carrés affectés d'un coefficient positif / négatif dans la décomposition de Gauss. Comme $1+a^2>0$ pour tout $a$, seul le signe de $a$ intervient.

\begin{itemize}
\item \underline{$a>0$} : trois coefficients strictement positifs. $\operatorname{sign}(q_a)=(3,0)$, $\operatorname{rg}(q_a)=3$, $\ker q_a=\{0\}$.
\item \underline{$a<0$} : deux coefficients positifs et un négatif. $\operatorname{sign}(q_a)=(2,1)$, $\operatorname{rg}(q_a)=3$, $\ker q_a=\{0\}$.
\item \underline{$a=0$} : $q_0(X)=(x_1+x_2)^2+x_3^2$. $\operatorname{sign}(q_0)=(2,0)$, $\operatorname{rg}(q_0)=2$. Le noyau est donné par $AX=0$ avec $A=\begin{pmatrix}1&1&0\\1&1&0\\0&0&1\end{pmatrix}$ : $x_1+x_2=0$ et $x_3=0$, d'où
\[
\ker q_0=\operatorname{Vect}\big((-1,1,0)\big).
\]
\end{itemize}
(Contrôle : $\det A=(1+a+a^2)\cdot a-a^2\cdot1=a(1+a^2)$, qui est nul si et seulement si $a=0$.)

\textbf{5. Base $q_a$-orthogonale.}

On cherche la base $\mathcal{B}=(e_1',e_2',e_3')$ dont les formes coordonnées sont $\ell_1,\ell_2,\ell_3$. Posons
\[
\begin{pmatrix}\ell_1\\ \ell_2\\ \ell_3\end{pmatrix}=Q\begin{pmatrix}x_1\\x_2\\x_3\end{pmatrix},\qquad Q=\begin{pmatrix}1&1&0\\0&1&-1\\0&0&1\end{pmatrix},\quad\det Q=1 .
\]
En résolvant : $x_3=\ell_3$, $x_2=\ell_2+\ell_3$, $x_1=\ell_1-\ell_2-\ell_3$, donc
\[
Q^{-1}=\begin{pmatrix}1&-1&-1\\0&1&1\\0&0&1\end{pmatrix}.
\]
Les colonnes de $Q^{-1}$ forment la base cherchée :
\[
\mathcal{B}=\left\{e_1'=\begin{pmatrix}1\\0\\0\end{pmatrix},\ e_2'=\begin{pmatrix}-1\\1\\0\end{pmatrix},\ e_3'=\begin{pmatrix}-1\\1\\1\end{pmatrix}\right\}.
\]
Dans cette base, $q_a=\ell_1^2+a\ell_2^2+(1+a^2)\ell_3^2$ n'a que des termes carrés : la matrice de $q_a$ dans $\mathcal{B}$ est $\operatorname{diag}(1,\,a,\,1+a^2)$, donc $\mathcal{B}$ est $q_a$-orthogonale (pour tout $a$, en particulier pour $a=2$).

Vérification directe, par exemple : $\varphi_a(e_1',e_2')={}^te_1'Ae_2'=(1,1,0)\cdot(-1,1,0)=0$ et $\varphi_a(e_2',e_3')=(0,a,-a)\cdot(-1,1,1)=0$.

\textbf{6. $q_a$ définie positive.}

$q_a$ est définie positive si et seulement si sa signature est $(3,0)$, c'est-à-dire si et seulement si $a>0$ (le coefficient $1+a^2$ est toujours strictement positif).
\[
\boxed{q_a \text{ est définie positive} \iff a\in\,]0,+\infty[}.
\]
(On retrouve ce résultat avec le critère de Sylvester : mineurs principaux $1>0$, $\begin{vmatrix}1&1\\1&1+a\end{vmatrix}=a>0$ et $\det A=a(1+a^2)>0$.)

\medskip\textbf{Exercice 3}

\textbf{1. $I=J$ sans calcul.}

Effectuons le changement de variable $x=\frac\pi2-u$ ($dx=-du$) dans $I$. Comme $\cos(\frac\pi2-u)=\sin u$ et $\sin(\frac\pi2-u)=\cos u$ :
\[
I=\int_{\pi/2}^{0}\frac{\sin^2u}{\sin u+\cos u}\,(-du)=\int_0^{\pi/2}\frac{\sin^2u}{\cos u+\sin u}\,du=J .
\]
(Autre méthode : $I-J=\displaystyle\int_0^{\pi/2}\frac{(\cos x-\sin x)(\cos x+\sin x)}{\cos x+\sin x}dx=\big[\sin x+\cos x\big]_0^{\pi/2}=1-1=0$.)

\textbf{2. $\cos x+\sin x=\sqrt2\cos\big(x-\frac\pi4\big)$.}

Par la formule $\cos(p-q)=\cos p\cos q+\sin p\sin q$ :
\[
\sqrt2\cos\Big(x-\frac\pi4\Big)=\sqrt2\Big(\frac{\sqrt2}{2}\cos x+\frac{\sqrt2}{2}\sin x\Big)=\cos x+\sin x .
\]
(L'énoncé original écrivait $2\cos(x-\frac\pi4)$ : c'est une coquille.) En particulier, $\cos x+\sin x>0$ sur $[0,\frac\pi2]$, donc $I$ et $J$ sont bien définies.

\textbf{3. Expression de $I+J$.}

Comme $\cos^2x+\sin^2x=1$ :
\[
I+J=\int_0^{\pi/2}\frac{dx}{\cos x+\sin x}=\int_0^{\pi/2}\frac{dx}{\sqrt2\cos(x-\frac\pi4)} .
\]
Avec $u=x-\frac\pi4$ ($du=dx$ ; $x=0\Rightarrow u=-\frac\pi4$ ; $x=\frac\pi2\Rightarrow u=\frac\pi4$) :
\[
I+J=\frac1{\sqrt2}\int_{-\pi/4}^{\pi/4}\frac{du}{\cos u}=\frac{\sqrt2}{2}\int_{-\pi/4}^{\pi/4}\frac{dx}{\cos x}.
\]

\textbf{4. Calcul de $I+J$, puis de $I$ et $J$.}

On pose $u=\tan\dfrac x2$ ; alors $dx=\dfrac{2\,du}{1+u^2}$ et $\cos x=\dfrac{1-u^2}{1+u^2}$, d'où $\dfrac{dx}{\cos x}=\dfrac{2\,du}{1-u^2}$.

\underline{Bornes.} Il faut $\tan\frac\pi8$. Avec $\tan\theta=\dfrac{2\tan\frac\theta2}{1-\tan^2\frac\theta2}$ et $\theta=\frac\pi4$ : en notant $T=\tan\frac\pi8>0$, $1=\dfrac{2T}{1-T^2}$, soit $T^2+2T-1=0$, d'où $T=-1+\sqrt2$. Ainsi $x=\pm\frac\pi4\Rightarrow u=\pm(\sqrt2-1)$.

\underline{Calcul.} Avec la décomposition $\dfrac{2}{1-u^2}=\dfrac{1}{1-u}+\dfrac{1}{1+u}$ :
\[
\int_{-\pi/4}^{\pi/4}\frac{dx}{\cos x}=\int_{1-\sqrt2}^{\sqrt2-1}\Big(\frac1{1+u}+\frac1{1-u}\Big)du=\Big[\ln\frac{1+u}{1-u}\Big]_{1-\sqrt2}^{\sqrt2-1}.
\]
Or, pour $u=\sqrt2-1$ : $\dfrac{1+u}{1-u}=\dfrac{\sqrt2}{2-\sqrt2}=\dfrac{\sqrt2(2+\sqrt2)}{2}=\sqrt2+1$, et pour $u=1-\sqrt2$ on obtient l'inverse $\dfrac1{\sqrt2+1}$. Donc
\[
\int_{-\pi/4}^{\pi/4}\frac{dx}{\cos x}=\ln(\sqrt2+1)-\ln\frac1{\sqrt2+1}=2\ln(1+\sqrt2),
\]
\[
\boxed{I+J=\frac{\sqrt2}{2}\times2\ln(1+\sqrt2)=\sqrt2\,\ln(1+\sqrt2)}.
\]
\remarque{On peut aussi utiliser directement une primitive de $\dfrac1{\cos x}$ sur $]-\frac\pi2,\frac\pi2[$ : $x\mapsto\ln\Big|\tan\Big(\dfrac x2+\dfrac\pi4\Big)\Big|$. On obtient $\ln\tan\frac{3\pi}8-\ln\tan\frac\pi8=\ln\frac{\sqrt2+1}{\sqrt2-1}=2\ln(1+\sqrt2)$.}

Comme $I=J$ et $I+J=\sqrt2\ln(1+\sqrt2)$ :
\[
\boxed{I=J=\frac{\sqrt2}{2}\ln(1+\sqrt2)\approx0{,}623}.
\]

\medskip\textbf{Exercice 4}

\textbf{1. Constantes $a$ et $b$.} Pour $|x|>1$, en réduisant au même dénominateur :
\[
\frac{a}{\sqrt{x^2-1}}+\frac{\sqrt{x^2-1}}{b}=\frac{ab+(x^2-1)}{b\sqrt{x^2-1}}=\frac{\frac1b\,x^2+\big(a-\frac1b\big)}{\sqrt{x^2-1}} .
\]
Par identification avec $\dfrac{x^2}{\sqrt{x^2-1}}$ : $\dfrac1b=1$ et $a-\dfrac1b=0$, soit $\boxed{a=1,\ b=1}$ :
\[
\frac{x^2}{\sqrt{x^2-1}}=\frac{1}{\sqrt{x^2-1}}+\sqrt{x^2-1}\qquad(\text{car } x^2=1+(x^2-1)).
\]

\textbf{2. Aire.}

Sur $[\sqrt2,2]$, $y=\dfrac{x^3}{\sqrt{x^2-1}}>0$, donc l'aire (en unités d'aire) est
\[
\mathcal{A}=\int_{\sqrt2}^{2}\frac{x^3}{\sqrt{x^2-1}}\,dx=\int_{\sqrt2}^{2}x\cdot\frac{x^2}{\sqrt{x^2-1}}\,dx=\int_{\sqrt2}^{2}\frac{x}{\sqrt{x^2-1}}\,dx+\int_{\sqrt2}^{2}x\sqrt{x^2-1}\,dx .
\]
On pose $u=x^2-1$, $du=2x\,dx$ ; $x=\sqrt2\Rightarrow u=1$ ; $x=2\Rightarrow u=3$ :
\[
\mathcal{A}=\frac12\int_1^3u^{-1/2}\,du+\frac12\int_1^3u^{1/2}\,du=\Big[\sqrt u\Big]_1^3+\frac12\Big[\frac23u^{3/2}\Big]_1^3=(\sqrt3-1)+\frac13(3\sqrt3-1).
\]
\[
\boxed{\mathcal{A}=2\sqrt3-\frac43\ \text{u.a.}\approx2{,}131\ \text{u.a.}}
\]

\end{spacing}
\end{document}
